Let $(x_k)_{k\ge 1}$ be an infinite sequence of real numbers in $[-1,1]$ such that for each $n$ the initial segment $x_1,\dots,x_n$ consists of distinct points (otherwise the Lagrange denominators vanish).
For each $n\ge 1$, define $l_k^{(n)}(x)$ (for $1\le k\le n$) to be the Lagrange fundamental polynomial associated to the nodes $x_1,\dots,x_n$:
\[
 l_k^{(n)}(x)=\frac{\prod_{1\le i\le n,\ i\ne k}(x-x_i)}{\prod_{1\le i\le n,\ i\ne k}(x_k-x_i)}.
\]
Define
\[
L_n(x):=\sum_{k=1}^n \bigl|l_k^{(n)}(x)\bigr|.
\]
Questions:
\begin{enumerate}
\item Must there exist $x\in(-1,1)$ such that $L_n(x) > \tfrac{2}{\pi}\log n - O(1)$ for infinitely many $n$?
\item Is it true that for almost all $x\in(-1,1)$,
\[\limsup_{n\to\infty} \frac{L_n(x)}{\log n}\ge \frac{2}{\pi}?\]
\end{enumerate}
